\begin{proof}[Proof of \cref{obj:lem:nuisance-frontier-converse}]
\begin{enumerate}
\item \textit{Public constants.}
Take the constants supplied by \cref{obj:thm:marked-component-certificate}, writing its amplitude constant as \(c_{\mathrm{cert}}\):
\[
0<c_{\mathrm{cert}}\le1,\qquad c_0>0,\qquad C>0.
\]
They depend only on the fixed public parameters. The parameter conditions give
\[
\gamma>0,\qquad S>0,\qquad \Delta>0.
\]
Define
\[
c_h=\frac14,\qquad
c_\delta=\min\{c_h,c_0\},\qquad
\rho_0=\min\left\{
c_{\mathrm{cert}},1,c_{\mathrm{cert}}c_h^\gamma,
\frac1{16C}
\right\}.
\]
Thus
\[
0<c_\delta\le c_h\le\frac12,\qquad
c_\delta\le c_0,\qquad
0<\rho_0\le1,
\]
and
\[
\rho_0\le c_{\mathrm{cert}},\qquad
\rho_0\le c_{\mathrm{cert}}c_h^\gamma,\qquad
C\rho_0\le\frac1{16}.
\]
Since \(d\ge1\), \(S>0\), and \(c_\delta,c_h\le1\),
\[
c_\delta^d\le c_\delta\le c_0,\qquad
c_\delta^S\le1,\qquad
c_h^d\le1,\qquad
c_\delta^{d+2S}\le1.
\]
Also,
\[
\rho_0^2\le\rho_0,\qquad \rho_0^4\le\rho_0,
\]
so the information constant satisfies
\[
C\rho_0^4c_h^dc_\delta^{d+2S}
\le C\rho_0\le\frac1{16}.
\]
Finally, define the separation constant and the claimed lower-bound constant by
\[
c_{\mathrm{sep}}=\rho_0^2c_\delta^S>0,
\qquad
c=\frac34c_{\mathrm{sep}}>0.
\]
% lean: CausalSmith.Stat.TwosamplePointcateAnnotationFrontier.nuisance_scaled_marked_priors

\item \textit{Admissible scales.}
Fix arbitrary \(n,m\in\mathbb N\) in the stated range, and write
\[
N=n+m,\qquad n\ge2,\qquad n\le N\le n^{q_*}.
\]
In particular, \(n,N>0\) and \(nN\ge1\). Choose
\[
h=c_h(nN)^{-1/\Delta},\qquad
\delta=c_\delta(nN)^{-\gamma/(S\Delta)},\qquad
a=\rho_0\delta^\alpha,\qquad
b=\rho_0\delta^\beta.
\]
The frontier condition implies
\[
S<S_{\mathrm{crit}}
=\frac{\gamma d}{2\gamma+d}<\gamma,
\]
because \(2\gamma+d>d>0\). Consequently,
\[
-\frac{\gamma}{S\Delta}\le-\frac1\Delta\le0.
\]
Monotonicity of positive-base powers for \(nN\ge1\) now gives
\[
0<\delta
\le c_\delta(nN)^{-1/\Delta}
\le c_h(nN)^{-1/\Delta}
=h
\le c_h\le\frac12.
\]
The amplitude definitions then give \(a,b>0\).
% lean: CausalSmith.Stat.TwosamplePointcateAnnotationFrontier.nuisance_scales_admissible

\item \textit{Separation and contrast calibration.}
Using \(\delta>0\) and \(S=\alpha+\beta\), multiplication of positive-base powers yields
\[
\begin{aligned}
ab
&=\rho_0^2\delta^{\alpha+\beta}\\
&=\rho_0^2c_\delta^S
  (nN)^{-\gamma S/(S\Delta)}\\
&=c_{\mathrm{sep}}(nN)^{-\gamma/\Delta}.
\end{aligned}
\]
Likewise,
\[
h^\gamma=c_h^\gamma(nN)^{-\gamma/\Delta}.
\]
The bounds \(\rho_0^2\le\rho_0\), \(c_\delta^S\le1\), and \(\rho_0\le c_{\mathrm{cert}}c_h^\gamma\) therefore imply
\[
\begin{aligned}
ab
&=\rho_0^2c_\delta^S(nN)^{-\gamma/\Delta}\\
&\le\rho_0(nN)^{-\gamma/\Delta}\\
&\le c_{\mathrm{cert}}c_h^\gamma(nN)^{-\gamma/\Delta}
=c_{\mathrm{cert}}h^\gamma.
\end{aligned}
\]
% lean: CausalSmith.Stat.TwosamplePointcateAnnotationFrontier.nuisance_scaled_marked_priors

\item \textit{Occupancy, class membership, and mixture distance.}
The definitions of the rate parameters give
\[
q_*=\frac{S_{\mathrm{crit}}}{S}\ge0,\qquad
\Delta=(2\gamma+d)(1+q_*),\qquad
\frac{\gamma d}{S\Delta}=\frac{q_*}{1+q_*}.
\]
Substituting the scale \(\delta\) and collecting powers therefore gives
\[
\begin{aligned}
N\delta^d
&=c_\delta^d
  n^{-q_*/(1+q_*)}N^{1/(1+q_*)}\\
&=c_\delta^d
  \left(\frac{N}{n^{q_*}}\right)^{1/(1+q_*)}
\le c_\delta^d\le c_0.
\end{aligned}
\]
Here the boundary condition makes the positive ratio \(N/n^{q_*}\) at most one, and \(1/(1+q_*)>0\). The remaining amplitude hypotheses follow from \(\rho_0\le c_{\mathrm{cert}}\):
\[
0<a\le c_{\mathrm{cert}}\delta^\alpha,\qquad
0<b\le c_{\mathrm{cert}}\delta^\beta.
\]

All hypotheses of \cref{obj:thm:marked-component-certificate} are thus satisfied for the construction of \cref{obj:def:marked-handle}. For its finite lattice index set \(\mathcal I\), write
\[
\mathcal S_{\mathrm{sign}}
=\{-1,+1\}^{\mathcal I}\times\{-1,+1\}^{\mathcal I}.
\]
The certificate supplies probability priors \(\varpi_\theta\) and, for every
\(\theta\in\{-1,+1\}\), \((\lambda,\eta)\in\mathcal S_{\mathrm{sign}}\), and \(x\in[0,1]^d\),
\[
P_{\theta,\lambda,\eta}\in\mathcal M,\qquad
\tau_{P_{\theta,\lambda,\eta}}(x)=\tau_\theta(x).
\]
At the evaluation profile this gives
\[
\tau_{P_{-1,\lambda,\eta}}(x_0)=2ab,\qquad
\tau_{P_{+1,\lambda,\eta}}(x_0)=-2ab.
\]

Define the mixtures of the original data blocks by
\[
M_\theta
=\int \mathcal L_P(\mathcal D_{n,m})\,d\varpi_\theta(P),
\qquad \theta\in\{-1,+1\}.
\]
The certificate also supplies
\[
H^2(M_{+1},M_{-1})\le CnNh^d\delta^da^2b^2.
\]
To evaluate its right-hand side, first note that
\[
\delta^da^2b^2=\rho_0^4\delta^{d+2S},
\qquad
d+\frac{\gamma(d+2S)}{S}
=d+\frac{\gamma d}{S}+2\gamma=\Delta.
\]
Hence
\[
\begin{aligned}
nNh^d\delta^da^2b^2
&=\rho_0^4c_h^dc_\delta^{d+2S}
  (nN)^{\,1-d/\Delta-\gamma(d+2S)/(S\Delta)}\\
&=\rho_0^4c_h^dc_\delta^{d+2S}.
\end{aligned}
\]
The inequality \(C\rho_0^4c_h^dc_\delta^{d+2S}\le1/16\) proves
\[
H^2(M_{+1},M_{-1})
\le C\rho_0^4c_h^dc_\delta^{d+2S}
\le\frac1{16}.
\]
% lean: CausalSmith.Stat.TwosamplePointcateAnnotationFrontier.nuisance_scaled_marked_priors

\item \textit{The target as a functional of the experiment law.}
The randomized observation space and the marked experiment model are
\[
\Omega_{\mathrm{obs}}
=
\bigl([0,1]^d\times\{0,1\}^2\bigr)^n
\times
\bigl([0,1]^d\times\{0,1\}\bigr)^m
\times[0,1],
\]
\[
\mathcal P_{\mathrm{exp}}
=
\left\{
\mathsf E_{n,m}(P_{\theta,\lambda,\eta}):
\theta\in\{-1,+1\},\
(\lambda,\eta)\in\mathcal S_{\mathrm{sign}}
\right\}.
\]
When \(m=0\), the auxiliary-record factor is the one-point empty product.

We show that equal marked experiment laws have equal point targets. Take two constructed laws
\[
P=P_{\theta,\lambda,\eta},\qquad
P'=P_{\theta',\lambda',\eta'},
\]
where both hypothesis indices belong to \(\{-1,+1\}\) and both sign pairs belong to \(\mathcal S_{\mathrm{sign}}\), and suppose
\[
\mathsf E_{n,m}(P)=\mathsf E_{n,m}(P').
\]
Since \(n>0\), projection onto the first labeled record gives
\[
\mathcal L_P(X,A,Y)=\mathcal L_{P'}(X,A,Y).
\]
Both covariate laws are uniform. Uniqueness of conditional probability kernels therefore gives
\[
\mathcal L_P(A,Y\mid X=x)
=
\mathcal L_{P'}(A,Y\mid X=x)
\quad\text{for Lebesgue-almost every }x\in[0,1]^d.
\]
For the independent Bernoulli construction, with the response selection of \cref{obj:synth_5}, the four conditional cells are
\[
\begin{aligned}
\Pr_P(A=1,Y=1\mid X=x)
 &=e_P(x)\mu_{1,P}(x),\\
\Pr_P(A=1,Y=0\mid X=x)
 &=e_P(x)\bigl(1-\mu_{1,P}(x)\bigr),\\
\Pr_P(A=0,Y=1\mid X=x)
 &=\bigl(1-e_P(x)\bigr)\mu_{0,P}(x),\\
\Pr_P(A=0,Y=0\mid X=x)
 &=\bigl(1-e_P(x)\bigr)\bigl(1-\mu_{0,P}(x)\bigr).
\end{aligned}
\]
The same formulas hold for \(P'\). Summing the first two equal cells identifies the propensity. The equal cells with \(Y=1\), together with overlap, then identify the arm means:
\[
e_P(x)=e_{P'}(x),\qquad
\mu_{0,P}(x)=\mu_{0,P'}(x),\qquad
\mu_{1,P}(x)=\mu_{1,P'}(x)
\quad\text{almost everywhere}.
\]
Indeed, the factors used for division satisfy
\[
e_P(x)\ge\varepsilon>0,\qquad 1-e_P(x)\ge\varepsilon>0.
\]
Conditional independence in the prescribed construction now gives
\[
\begin{aligned}
\mathcal L_P(A,Y(0),Y(1)\mid X=x)
&=\operatorname{Bernoulli}(e_P(x))
  \otimes\operatorname{Bernoulli}(\mu_{0,P}(x))
  \otimes\operatorname{Bernoulli}(\mu_{1,P}(x))\\
&=\mathcal L_{P'}(A,Y(0),Y(1)\mid X=x)
\quad\text{almost everywhere}.
\end{aligned}
\]
Integration against their common uniform covariate law proves equality of the full primitive laws.

Their common primitive law thus gives identical Bernoulli margin laws for \(A,Y(0),Y(1)\), so admissible conditional-margin versions agree almost everywhere. The positive Hölder orders give continuity of \(e_P,\mu_{0,P},\tau_P\), and
\[
\mu_{1,P}=\mu_{0,P}+\tau_P
\quad\text{on }[0,1]^d
\]
gives continuity of the treated mean; the same holds for \(P'\). Since the uniform law has full support on the cube, these continuous versions agree everywhere. In particular,
\[
\tau_P(x_0)=\tau_{P'}(x_0).
\]

We may consequently define a functional on all nonnegative Borel measures on \(\Omega_{\mathrm{obs}}\) by
\[
\Psi(\mathbb P)=
\begin{cases}
\tau_{P_{\theta,\lambda,\eta}}(x_0),
&
\mathbb P=\mathsf E_{n,m}(P_{\theta,\lambda,\eta})
\text{ for some }
\theta\in\{-1,+1\},\
(\lambda,\eta)\in\mathcal S_{\mathrm{sign}},\\
0,&\mathbb P\notin\mathcal P_{\mathrm{exp}}.
\end{cases}
\]
The equality just proved makes the first case independent of the representative. Thus
\[
\Psi\bigl(\mathsf E_{n,m}(P_{\theta,\lambda,\eta})\bigr)
=\tau_{P_{\theta,\lambda,\eta}}(x_0)
\]
throughout the marked family.
% lean: CausalSmith.Stat.TwosamplePointcateAnnotationFrontier.marked_experiment_target_exists

\item \textit{A common prior for testing.}
Use the finite index space
\[
\mathcal Z_{\mathrm{test}}
=\mathcal S_{\mathrm{sign}}\times\mathcal S_{\mathrm{sign}},
\qquad
z_{\mathrm{test}}
=((\lambda_-,\eta_-),(\lambda_+,\eta_+)).
\]
Define the common prior by its point masses:
\[
\omega(\{z_{\mathrm{test}}\})
=
w_{-1}(\lambda_-,\eta_-)\,
w_{+1}(\lambda_+,\eta_+).
\]
The nonnegative weights and their unit sums from \cref{obj:thm:marked-component-certificate} make \(\omega\) a probability measure. Define the two probability kernels by
\[
\mathbb B_{z_{\mathrm{test}}}
=\mathsf E_{n,m}(P_{-1,\lambda_-,\eta_-}),
\qquad
\mathbb C_{z_{\mathrm{test}}}
=\mathsf E_{n,m}(P_{+1,\lambda_+,\eta_+}).
\]
Their values belong to \(\mathcal P_{\mathrm{exp}}\). For every
\(z_{\mathrm{test}},z'_{\mathrm{test}}\in\mathcal Z_{\mathrm{test}}\), the target identities give
\[
\Psi(\mathbb B_{z_{\mathrm{test}}})
-\Psi(\mathbb C_{z'_{\mathrm{test}}})
=2ab-(-2ab)=4ab>0.
\]

Summing out the unused sign pair in each kernel gives the mixture identities
\[
\int\mathbb B_{z_{\mathrm{test}}}\,d\omega(z_{\mathrm{test}})
=M_{-1}\otimes\operatorname{Unif}([0,1]),
\]
\[
\int\mathbb C_{z_{\mathrm{test}}}\,d\omega(z_{\mathrm{test}})
=M_{+1}\otimes\operatorname{Unif}([0,1]).
\]
The uniform factor is the independent randomizer. By the common-factor identity in \cref{obj:lem:hellinger-block-calculus} and symmetry of the squared difference in \cref{obj:def:hellinger},
\[
\begin{aligned}
\mathsf h^2\left(
\int\mathbb B_{z_{\mathrm{test}}}\,d\omega,
\int\mathbb C_{z_{\mathrm{test}}}\,d\omega
\right)
&=H^2(M_{-1},M_{+1})\\
&=H^2(M_{+1},M_{-1})
\le\frac1{16}.
\end{aligned}
\]
% lean: CausalSmith.Stat.TwosamplePointcateAnnotationFrontier.marked_common_prior_hellinger

\item \textit{Testing implies a minimax bound.}
Apply \cref{obj:lem:published-fuzzy-testing} to the model \(\mathcal P_{\mathrm{exp}}\), the functional \(\Psi\), and the preceding kernels and prior, with
\[
v_0=1,\qquad s_0=4ab>0,\qquad
u_0=\frac1{16}\in[0,2).
\]
Here one observation is the entire randomized data block; the one-fold product space is identified with that block by its sole coordinate. For every Borel real-valued decision \(T\) on \(\Omega_{\mathrm{obs}}\), the testing bound gives
\[
\sup_{\mathbb P\in\mathcal P_{\mathrm{exp}}}
\int_{\Omega_{\mathrm{obs}}}
\left|T(w_{\mathrm{obs}})-\Psi(\mathbb P)\right|
\,d\mathbb P(w_{\mathrm{obs}})
\ge
ab\left(1-\frac{\sqrt{63}}{32}\right).
\]
The numerical overlap factor satisfies
\[
63\le64,\qquad
\frac{\sqrt{63}}{32}\le\frac14,\qquad
1-\frac{\sqrt{63}}{32}\ge\frac34.
\]
Moreover, the target representation and \cref{obj:def:risk} give
\[
\begin{aligned}
&\sup_{\mathbb P\in\mathcal P_{\mathrm{exp}}}
\int_{\Omega_{\mathrm{obs}}}
\left|T(w_{\mathrm{obs}})-\Psi(\mathbb P)\right|
\,d\mathbb P(w_{\mathrm{obs}})\\
&\qquad=
\sup_{\substack{\theta\in\{-1,+1\}\\
(\lambda,\eta)\in\mathcal S_{\mathrm{sign}}}}
\mathcal R_{n,m}(T,P_{\theta,\lambda,\eta})
\le
\sup_{P\in\mathcal M}\mathcal R_{n,m}(T,P),
\end{aligned}
\]
where the last inequality uses the certified class membership. Thus every such decision satisfies
\[
\sup_{P\in\mathcal M}\mathcal R_{n,m}(T,P)\ge\frac34ab.
\]
Taking the infimum over the decisions in \cref{obj:def:minimax} proves
\[
R(n,m)\ge\frac34ab.
\]
All loss integrals and suprema here retain their extended nonnegative interpretation.
% lean: CausalSmith.Stat.TwosamplePointcateAnnotationFrontier.marked_fuzzy_testing_minimax

\item \textit{Substitution of the separation scale.}
The identity \(ab=c_{\mathrm{sep}}(nN)^{-\gamma/\Delta}\) now yields
\[
R(n,m)
\ge\frac34ab
=\frac34c_{\mathrm{sep}}(nN)^{-\gamma/\Delta}
=c(nN)^{-\gamma/\Delta}.
\]
The positive constant \(c\) depends only on the fixed public parameters, and the sample sizes were arbitrary in the stated range.
% lean: CausalSmith.Stat.TwosamplePointcateAnnotationFrontier.nuisance_frontier_converse
\end{enumerate}
\end{proof}
